% !TeX program = pdflatex
% Self-contained source: pdflatex twice; no fontspec or external bibliography.
\documentclass[11pt,a4paper]{article}
\usepackage[margin=24mm,headheight=14pt]{geometry}
\usepackage[T1]{fontenc}
\usepackage[utf8]{inputenc}
\usepackage{lmodern,microtype}
\usepackage{amsmath,amssymb,mathtools,amsthm}
\usepackage{enumitem,booktabs,longtable,array,xurl,listings,needspace}
\usepackage[hidelinks,unicode]{hyperref}
\hypersetup{pdftitle={Economic Theory: Proof Attempts and Adversarial Checks},pdfauthor={Research notes},pdfsubject={Eleven user-supplied mathematical problems}}
\newcommand{\NE}{\operatorname{NE}}
\DeclareMathOperator*{\argmax}{arg\,max}
\DeclareMathOperator*{\argmin}{arg\,min}
\newtheorem{theorem}{Theorem}[section]
\newtheorem{lemma}[theorem]{Lemma}
\newtheorem{proposition}[theorem]{Proposition}
\newtheorem{corollary}[theorem]{Corollary}
\theoremstyle{definition}
\newtheorem{definition}[theorem]{Definition}
\theoremstyle{remark}
\newtheorem*{remark}{Remark}
\setlength{\parindent}{0pt}
\setlength{\parskip}{5pt plus 1pt minus 1pt}
\setlength{\emergencystretch}{2em}
\setcounter{tocdepth}{1}
\setlist{leftmargin=*,itemsep=2pt,topsep=4pt}
\lstset{basicstyle=\small\ttfamily,columns=fullflexible,keepspaces=true,
 breaklines=true,breakatwhitespace=false,showstringspaces=false,
 frame=single,framerule=0.25pt,aboveskip=8pt,belowskip=8pt,xleftmargin=3pt,xrightmargin=3pt}
\allowdisplaybreaks[1]
\raggedbottom

\begin{document}
{\large\bfseries Research attempt and adversarial verification}\par
9 October 2026\hfill Original-target status: \textbf{UNRESOLVED}\par
\section{C73-11: Universal convergence of simultaneous fictitious play}
\label{sec:C73-11}
\noindent\textbf{Input:} \texttt{C73-11.tex}. \quad\textbf{Tools:} web browsing [YES]; Python computation [YES].

\subsection*{1) FORMAL RESTATEMENT}
Let $N=\{1,\ldots,n\}$, with $n\geq1$, and let every finite action set $S_i$ be
nonempty. Let $X=\prod_i\Delta(S_i)$, and extend each payoff array $u_i$
multilinearly to independent mixed profiles. For every $x^1\in X$, and every
sequence of legal simultaneous best replies, set
\[
 a_i^{t+1}\in\argmax_{a\in S_i}u_i(a,x_{-i}^t),\qquad
 x_i^{t+1}=x_i^t+\frac{e_{a_i^{t+1}}-x_i^t}{t+1},\quad t\geq1.
\]
The class to characterize is
\[
 \mathcal C_{\rm FP}=\{G:\ \forall x^1\in X\ \forall\text{ legal reply sequences},
 \ \operatorname{dist}(x^t,\NE(G))\longrightarrow0\}.
\]
Euclidean distance is used. Boundary initializations and arbitrary choices at ties
are included. The output sought is a noncircular condition on payoff arrays,
not another formula quantifying over the same trajectories. This is a classification
agenda; a nonconvergent game excludes that game, not the classification question itself.

\subsection*{2) QUICK LITERATURE/CONTEXT CHECK}
The input survey \cite{li-survey} supplies the learning-classification context.
Monderer and Shapley \cite{monderer-fp} prove a fictitious-play property for
identical-interest games. The proof below independently treats positive weighted
potential games under the precise simultaneous update and tie convention in the
input. A specific nonconvergent $3\times3$ array is analyzed directly, rather than
relying on an unspecified version of a cyclic-play example.

\subsection*{3) ATTACK PLAN}
\textbf{Proof track.} Use a continuous Nash residual; seek a payoff-derived
Lyapunov function with a summable discretization error; turn weighted summability
of residuals into pointwise residual convergence.

\textbf{Construction track.} Test a cyclic $3\times3$ interaction, first with
synchronized beliefs and then with distinct pure initializations. Track constant
best-reply blocks exactly to distinguish an actual counterexample from a finite
simulation that only appears to cycle.

\textbf{Tools and purpose.} Multilinear extension identifies directional gains;
compactness equates residual and Nash-set convergence; Taylor's theorem controls
simultaneous-update errors; harmonic step sums prevent isolated residual spikes;
strict payoff comparisons make block lengths exact; permutation support arguments
identify the equilibrium set; integer count arithmetic eliminates floating-point
tie errors in the search.

\subsection*{4) WORK}
For later use, define
\[
 r_i(x)=\max_{a_i}u_i(a_i,x_{-i})-u_i(x),\qquad R(x)=\sum_i r_i(x).
\]
Every $r_i$ is nonnegative and Lipschitz on $X$, because it is a finite maximum
of polynomials minus a polynomial. Also $R(x)=0$ exactly on $\NE(G)$: no mixed
deviation can outperform the best pure deviation. Since $X$ is compact and
$\NE(G)$ is nonempty by finite-game Nash existence \cite{nash}, any sequence has
$R(x^t)\to0$ if and only if its distance to $\NE(G)$ tends to zero. To verify the
nontrivial direction, a sequence with distance bounded below has a convergent
subsequence whose limit is outside $\NE(G)$ but has zero residual, a contradiction.

\begin{theorem}[Positive weighted potential games belong to $\mathcal C_{\rm FP}$]
Suppose there are $w_i>0$ and an array $\Phi$ such that, for every pure opponents'
profile and every $a_i,b_i$,
\[
 u_i(a_i,s_{-i})-u_i(b_i,s_{-i})
 =w_i\bigl(\Phi(a_i,s_{-i})-\Phi(b_i,s_{-i})\bigr).
\]
Then every process in the formal restatement converges in distance to $\NE(G)$.
\end{theorem}
\begin{proof}
Extend $\Phi$ multilinearly and put $\alpha_t=1/(t+1)$ and
$d^t=(e_{a_i^{t+1}}-x_i^t)_i$. The potential identity extends by multilinearity,
so
\[
 \langle\nabla\Phi(x^t),d^t\rangle
 =\sum_i\frac{r_i(x^t)}{w_i}=:Q(x^t)\geq0.
\]
The norm of $d^t$ is bounded by a constant $D$, and the Hessian of $\Phi$ is bounded
on the compact convex set $X$. Taylor's theorem on the line segment of the update
therefore gives a constant $C$ with
\[
 \Phi(x^{t+1})-\Phi(x^t)\geq\alpha_tQ(x^t)-C\alpha_t^2.
\]
Writing $e_t=C\alpha_t^2$, the bounded sequence
$\Phi(x^t)-\sum_{s=t}^{\infty}e_s$ is nondecreasing. It converges, and summing the
displayed inequality gives $\sum_t\alpha_tQ(x^t)<\infty$.

Here is the required no-spike argument. Let $L$ be a Lipschitz constant for $Q$.
If $LD=0$, then $Q(x^t)$ is constant along the process and the divergent sum
$\sum_t\alpha_t$ forces that constant to be zero. Otherwise suppose infinitely
many $t$ have $Q(x^t)\geq\delta>0$. Starting at any sufficiently large such time,
let $k>t$ be the first index such that
$\sum_{s=t}^{k-1}\alpha_s\geq\delta/(2LD)$. For $t\leq j<k$, the minimality of
$k$ gives $Q(x^j)\geq\delta/2$. Thus this interval contributes at least
$\delta^2/(4LD)$ to $\sum_s\alpha_sQ(x^s)$. If there are infinitely many spikes,
choose these finite intervals successively with disjoint interiors. Their
contributions contradict the finite sum. Hence $Q(x^t)\to0$. Positive weights
make this equivalent to $R(x^t)\to0$, and the residual criterion finishes the proof.
\end{proof}
This is a sufficient array condition, not a necessary one. For example, strict
dominance gives another class: if every player has a strictly dominant pure action,
all selected replies are those actions, and
$x_i^t=e_{a_i^*}+(x_i^1-e_{a_i^*})/t$. The arrays
\[
 A=\begin{pmatrix}3&2\\1&0\end{pmatrix},\qquad
 B=\begin{pmatrix}3&0\\2&0\end{pmatrix}
\]
have strictly dominant first actions but cannot have a positive weighted potential:
the row player's rectangular payoff difference is zero, while the column player's
is one; multiplying either by a positive weight cannot reconcile them.

\begin{theorem}[An explicit game outside $\mathcal C_{\rm FP}$]
Let
\[
 A=\begin{pmatrix}0&0&1\\1&0&0\\0&1&0\end{pmatrix},\qquad
 B=\begin{pmatrix}0&1&0\\0&0&1\\1&0&0\end{pmatrix}.
\]
Start from $x^1=e_1$, $y^1=e_2$ and choose the lowest numbered best action at every
tie. This legal process does not approach the Nash set.
\end{theorem}
\begin{proof}
Each player best responds by selecting the cyclic successor of a largest
coordinate of the other's empirical counts. First, the only equilibrium is
$x=y=(1/3,1/3,1/3)$. Indeed, a singleton support for one player forces a singleton
support for the other; their cyclic-successor requirements would give a pure
action equal to its successor twice, which is impossible. If a support has size
two, the other player's largest-coordinate set must contain at least two actions.
The other support cannot have size one; if it had size three, best-response
indifference would force the first player to be uniform with full support.
Consequently both supports have size two and both distributions put weight
$1/2$ on their supports. Each support must be the cyclic successor of the other,
so a two-element subset would be invariant under a three-cycle squared. No such
subset exists. Full supports force both distributions to be uniform.

Use integer counts $c=tx^t$ and $d=ty^t$. These are integers for the chosen initial
condition. Consecutive constant reply blocks have the following forms; the final
column is the condition at which the next block begins:
\[
\begin{array}{c|c|c}
 k\bmod6&(\text{row action},\text{column action})&\text{switch condition}\\\hline
 0&(3,2)&c_3=c_1\\
 1&(3,1)&d_1=d_2\\
 2&(2,1)&c_2=c_3+1\\
 3&(2,3)&d_3=d_1\\
 4&(1,3)&c_1=c_2\\
 5&(1,2)&d_2=d_3+1
\end{array}
\]
We verify this table and its repetition, including all tie offsets. At the start
of a six-block cycle write $c=(a,b,c_0)$ and $d=(d_0,e,f)$, with
$a>b\geq c_0$, $e=f+1$, and $f\geq d_0$. Initially these are $(1,0,0)$ and
$(0,1,0)$. The six positive lengths are
\begin{align*}
 A_0&=a-c_0, &B_0&=e+A_0-d_0, &C_0&=a+B_0-b+1,\\
 D_0&=C_0+A_0+1, &E_0&=D_0+B_0+1, &F_0&=E_0+C_0+1.
\end{align*}
In the first block, the third row count catches the first; at equality the column's
first action takes precedence over its second. In the second, the first column
count catches the second; the row's second action takes precedence over its third.
In the third, the second row count must exceed the third, since equality still
favors the column's first action. In the fourth and fifth, equality suffices by
the same lowest-action rule; in the sixth the second column count must exceed
the third. These are exactly the six conditions in the table. The untouched
third competitor is never larger than the two being compared, by the displayed
initial inequalities and the positive preceding increments.

At the next cycle the counts are
\[
 c'=(a+E_0+F_0,\ b+C_0+D_0,\ a+B_0),\quad
 d'=(e+A_0+C_0,\ e+A_0+F_0,\ e+A_0+C_0+E_0).
\]
They satisfy the same inequalities, now strictly where needed, and $d'_2=d'_3+1$.
This proves repetition of the table. If $\ell_k$ is the length of block $k$, the
formulas also give
\[
 \ell_0=1,\quad\ell_1=2,\quad\ell_2=4,\qquad
 \ell_{k+3}=\ell_{k+2}+\ell_k+1\quad(k\geq0).
\]
For clarity, the cross-cycle identities follow by subtraction:
$A'_0=D_0+F_0+1$, $B'_0=A'_0+E_0+1$, and
$C'_0=B'_0+F_0+1$. Thus the recurrence covers every index, not just the
three within-cycle transitions.

The lengths increase and $\ell_{k+3}\geq2\ell_k$. Grouping the earlier lengths
in triples consequently gives $\sum_{j=0}^k\ell_j\leq6\ell_k$.
If block $k$ begins at time $t_k=1+\sum_{j<k}\ell_j$, its averaging weight is
\[
 \lambda_k=\frac{\ell_k}{t_k+\ell_k}\geq\frac17.
\]
During this block the row player repeats one action $a$, so
$x^{t_k+\ell_k}=(1-\lambda_k)x^{t_k}+\lambda_ke_a$.
If both endpoint profiles were within Euclidean distance $1/30$ of the uniform
profile, their $a$ coordinates could differ by at most $1/15$. But the update
would increase that coordinate by at least
\[
 \frac17\left(\frac23-\frac1{30}\right)=\frac{19}{210}>\frac1{15}.
\]
At least one endpoint is therefore a distance at least $1/30$ from the unique
Nash equilibrium. The endpoint times tend to infinity. Set convergence fails.
\end{proof}

\subsection*{5) VERIFICATION}
For a one-player game, every chosen action maximizes its fixed payoff and initial
mass on suboptimal actions is divided by $t$. The same argument handles singleton
action sets. In the two-player identical-payoff anti-coordination game with matrix
$\left(\begin{smallmatrix}0&1\\1&0\end{smallmatrix}\right)$, synchronized pure-start
fictitious play has total residual $(t+1)/t^2$ at odd $t$, consistent with the
potential theorem. Actual actions need not converge in this example.

The integer implementation of the cyclic counterexample records total residuals
$0.3054$, $0.282011$, $0.27850335$, and $0.278920207$ at times
$10^2,10^3,10^4,10^5$. The analytical block proof, not those four numbers,
certifies nonconvergence. An initial attempted test with both initial beliefs
$e_1$ instead produced residuals tending to zero; the universal quantifier requires
only one legal bad initialization, not that every initialization fail.

Minimal search pseudocode is
\begin{lstlisting}
c, d = (1,0,0), (0,1,0)
for t = 1,2,...:
    row = first_argmax(A*d)
    col = first_argmax(transpose(B)*c)
    record_switches_and_residual_before_update()
    c[row] += 1; d[col] += 1
\end{lstlisting}
The potential proof uses $\sum\alpha_t^2<\infty$ and
$\sum\alpha_t=\infty$ separately. It neither assumes strict potential improvement
at a tie nor interchanges a maximum and an infinite-time limit.

\Needspace{8\baselineskip}
\subsection*{6) FINAL}
\textbf{LABEL: UNRESOLVED.}

\textbf{Strongest checked partial result.} A complete convergence theorem for positive
weighted potential games, a distinct dominance class, and a fully verified explicit
$3\times3$ exclusion with an infinite subsequence bounded away from its unique equilibrium.

\textbf{Exact first gap.} No noncircular necessary-and-sufficient payoff-array
criterion for arbitrary finite games is established.

\textbf{Next three moves.}
\begin{enumerate}[nosep]
\item Characterize which nonpotential perturbations preserve the Lyapunov inequality
with an error that remains summable along every legal path.
\item Search for cyclic-block certificates with rational inequalities and prove
completeness or incompleteness of that certificate family in small dimensions.
\item Test proposed criteria against both the strictly dominant nonpotential array
and the synchronized versus asymmetric paths of the cyclic array.
\end{enumerate}
\textbf{Counterexample perspective.} The displayed game falsifies universal convergence
across all games. It does not falsify a yet-unspecified characterization, and it
does not show that all nonpotential games fail.

\paragraph{Reproducibility and scope.} This is one of eleven companion reports. The bundle includes \texttt{verify.py} and its executed results. Finite experiments supplement the proofs; they are not proofs of asymptotic claims. Literature coverage is targeted, and no novelty or formal proof-assistant certification is claimed.

\begin{thebibliography}{99}
\small
\bibitem{li-survey}
Hanyu Li, Wenhan Huang, Zhijian Duan, David Henry Mguni, Kun Shao, Jun Wang, and Xiaotie Deng.
\emph{A survey on algorithms for Nash equilibria in finite normal-form games}. arXiv:2312.11063 (2023).
\url{https://arxiv.org/abs/2312.11063}.

\bibitem{monderer-fp}
Dov Monderer and Lloyd S. Shapley.
Fictitious Play Property for Games with Identical Interests.
\emph{Journal of Economic Theory} 68(1) (1996), 258--265.
\url{https://www.sciencedirect.com/science/article/pii/S0022053196900149}.

\bibitem{nash}
John F. Nash, Jr. Equilibrium points in $n$-person games.
\emph{Proceedings of the National Academy of Sciences} 36(1) (1950), 48--49.
DOI: 10.1073/pnas.36.1.48.
\url{https://www.pnas.org/doi/10.1073/pnas.36.1.48}.

\end{thebibliography}

\end{document}
